\section{Ground states and phase transfer}\label{setup:section}

The probabilistic part of the proof uses the squared modulus of a ground
state. This section justifies that reduction and identifies the additional
estimate needed to recover the occupation of an arbitrary complex state.
The relevant phase information is local: moving one particle with the bath
fixed preserves its phase whenever the motion stays in the allowed region.

\subsection{Length units and the occupation fraction}

We work on the scaled torus $\Lambda_K=(\R/K\Z)^3$, with integer side
$K$, volume $V=K^3$, and exclusion distance $r>0$. Its distance is denoted
by $d$. For every integer $n\ge0$ put
\[
 \Omega_n=\{(x_1,\ldots,x_n)\in\Lambda_K^n:
                    d(x_i,x_j)>r\text{ for }i<j\},
\]
with $\Omega_0$ a single point. The full, labelled-particle Dirichlet form,
without a symmetry restriction, is
\[
 q_n[f]=\sum_{i=1}^n\int_{\Omega_n}|\nabla_i f|^2,
 \qquad D(q_n)=H_0^1(\Omega_n).
\]
Here $H_0^1$ is the closure of smooth functions compactly supported in the
open configuration set. The bosonic form is its restriction to functions
invariant under coordinate permutations. Write $E_n$ for the full bottom
energy when $\Omega_n$ is nonempty, set $E_n=+\infty$ otherwise, and
put $E_0=0$. We consider nonempty $\Omega_N$. The particle-insertion
argument in the next section proves nonemptiness for every additional
particle number at which an energy is used. Every wave function is
extended by zero outside its allowed set.

The mean particle count per unit cell and its product with the exclusion
distance are
\begin{equation}\label{setup:scales}
 D=\frac{N}{K^3},\qquad \alpha=Dr,
 \qquad \frac{\alpha_0}{2}\le\alpha\le2\alpha_0,
 \qquad D\ge D_0.
\end{equation}
The numerical constant $\alpha_0$ will be fixed sufficiently large, and
$D_0$ will subsequently be increased. Partition the torus into unit cells
$B_v=v+[-1/2,1/2)^3$, indexed by
$\mathcal L=(\Z/K\Z)^3$; these sets are understood modulo $K$.
All estimates are uniform in the ground state and in the displayed
$\alpha$ window. A constant may depend on parameters already fixed, but
not on $N$ or $K$. Volume thresholds may depend on the fixed density;
where needed, they will be uniform when $D$ ranges over a compact
subinterval of $[D_0,\infty)$.

To relate these units to the original torus of side $L$ and exclusion
$a$, choose a length unit $s=L/K$ and set $r=a/s$. The map
\[
 \widetilde\Psi(z_1,\ldots,z_N)
       =s^{3N/2}\Psi(sz_1,\ldots,sz_N)
\]
is unitary and multiplies the quadratic form by $s^2$. It therefore
preserves ground states. The gas parameter is unchanged:
\begin{equation}\label{setup:gas-parameter}
 Dr^3=\frac{N}{L^3}a^3=\frac{\alpha^3}{D^2}.
\end{equation}
Thus, with $\alpha$ in its fixed window, increasing the particle count
$D$ per unit cell corresponds to increasing physical dilution. The integer
$K$ will be chosen at the end of the proof.

Let $\Psi$ be a normalized, possibly complex, bosonic ground state on
$\Lambda_K$. Write $Y=(x_2,\ldots,x_N)$ and define
\[
 \Gamma^{(1)}_\Psi(x;y)
   =N\int\Psi(x,Y)\overline{\Psi(y,Y)}\,\dd Y,
 \qquad \varphi_0(x)=V^{-1/2}.
\]
The operator with kernel $\Psi(x,Y)$ is Hilbert--Schmidt with squared
norm one, so $\Gamma^{(1)}_\Psi$ is positive, trace class, and has trace
$N$. Consequently the normalized constant-orbital occupation is
\begin{equation}\label{setup:occupation}
 B(\Psi)=\frac{\langle\varphi_0,
                       \Gamma^{(1)}_\Psi\varphi_0\rangle}{N}
       =\frac1V\int\left|\int\Psi(x,Y)\,\dd x\right|^2\dd Y.
\end{equation}
Changing variables in this identity proves its invariance under the
unitary rescaling above. For the rest of the proof put
\begin{equation}\label{setup:position-law}
 \Phi=|\Psi|,\qquad \mu(\dd X)=\Phi(X)^2\,\dd X,
 \qquad m(Y)=\int\Phi(x,Y)^2\,\dd x.
\end{equation}
Thus $m(Y)\dd Y$ is the bath marginal of the probability measure $\mu$.

\subsection{The full ground equation and its equality case}

The componentwise positivity mechanism is classical; compare
Faris and Simon \cite[Lemmas~1 and~3]{FarisSimon1975GroundStates}.
We give the form argument to include the symmetry restriction and the
possibly disconnected torus domain. The multiplier identity below is the
classical ground-state representation; see
\cite[Section~2.3]{FrankSeiringer2008}. Its direct proof here supplies the
required form-domain statement.

\begin{lemma}[Ground states and component phases]\label{setup:ground}
For every integer $n\ge1$ with nonempty $\Omega_n$, the full and
bosonic Dirichlet operators
have compact resolvent and the same bottom energy $E_n$.
In particular, both $\Psi$ and $\Phi=|\Psi|$ are full ground
eigenfunctions at energy $E_N$. They have smooth representatives in
$\Omega_N$. On every connected component $C$ of $\Omega_N$, either
both vanish identically or
\[
 \Phi>0,\qquad \Psi=e^{i\vartheta_C}\Phi
\]
for one constant $\vartheta_C\in\R$.
For every real bounded Lipschitz function $u$ on $\Lambda_K^N$,
\begin{equation}\label{setup:multiplier}
 q_N[u\Phi]-E_N\norm{u\Phi}_2^2
                 =\int\Phi^2|\nabla u|^2.
\end{equation}
The gradient on the right includes all particle coordinates.
\end{lemma}

\begin{proof}
Zero extension maps $H_0^1(\Omega_n)$ isometrically into
$H^1(\Lambda_K^n)$. On a bounded set in that Sobolev space, the squared
$L^2$ Fourier tail outside frequencies of length at most $R$ is bounded
by $C R^{-2}$. The remaining Fourier space is finite dimensional. This
proves compactness of the form embedding, hence compact resolvent and
attainment of the bottom energy. Restricting to the closed symmetric
subspace preserves compactness.

For $f\in H_0^1(\Omega_n)$ define its root-mean-square symmetrization by
\[
 F(X)=\left(\frac1{n!}\sum_{\pi\in S_n}|f(\pi X)|^2\right)^{1/2}.
\]
The gradient inequality for the Euclidean norm of a vector gives
$\norm F_2=\norm f_2$ and $q_n[F]\le q_n[f]$. It also preserves
zero Dirichlet data. More explicitly, apply the construction to smooth
compactly supported approximants of $f$. Their vector norms have compact
support in $\Omega_n$ and admit interior smooth approximation. They
converge to $F$ in $L^2$ and remain bounded in $H^1$. Weak closedness of
$H_0^1$ and lower semicontinuity prove the assertion and the energy
inequality for $f$. Since $F$ is symmetric, the bosonic bottom is at
most the full bottom. Inclusion of the domains gives the reverse
inequality.

The bosonic minimizer $\Psi$ therefore minimizes the full Rayleigh
quotient. First variation against arbitrary full form-domain functions
gives its full weak eigenfunction equation. The modulus contraction and
the full variational principle yield
\[
 E_N\le q_N[\Phi]\le q_N[\Psi]=E_N,
 \qquad \norm\Phi_2=1.
\]
Thus $\Phi$ satisfies the same full weak eigenfunction equation.
Interior elliptic regularity makes both functions smooth in the open
set. Their almost-everywhere identity $\Phi=|\Psi|$ holds everywhere
there by continuity. Since $\Delta\Phi=-E_N\Phi\le0$, the strong
minimum principle makes $\Phi$ either strictly positive on a connected
component or identically zero there.

On a positive component set $v=\Psi/\Phi$. It is smooth and has
$|v|=1$, so $\operatorname{Re}(\overline v\nabla v)=0$ and
\[
 |\nabla\Psi|^2=|\nabla\Phi|^2+\Phi^2|\nabla v|^2.
\]
There are at most countably many open components. Summing their energy
identities and using $q_N[\Psi]=q_N[\Phi]$ shows that the last,
nonnegative term has zero integral on each. Hence $\nabla v=0$, and
$v$ is constant on that component. On a component where $\Phi=0$,
$\Psi$ is also zero. This leaves unrelated phases on distinct positive
components; no assertion of global connectedness or uniqueness is used.

Finally, multiplication by a bounded Lipschitz function preserves the
form domain. Test the weak equation for $\Phi$ with $u^2\Phi$ and
expand its gradient. Subtracting the resulting identity from the
expansion of $q_N[u\Phi]$ gives Equation~\eqref{setup:multiplier}.
\end{proof}

\begin{lemma}[Allowed slices]\label{setup:slices}
For almost every fixed bath $Y$, a function $f\in H_0^1(\Omega_N)$
restricts to an element of $H_0^1(S_Y)$, where
\[
 S_Y=\{x\in\Lambda_K:(x,Y)\in\Omega_N\}.
\]
The same assertion holds when any nonempty subset of coordinates is
left unfixed. For the smooth ground representatives, every connected
subset of $S_Y$ has one common phase on its nonzero points, or the wave
function vanishes throughout it.
\end{lemma}

\begin{proof}
Choose $f_j\in C_c^\infty(\Omega_N)$ converging to $f$ in the form
norm, and zero-extend. Fubini's theorem gives a subsequence converging
in $H^1$ in the unfixed coordinates for almost every bath. For each
$j$, the restricted function is smooth and compactly supported in its
allowed slice, because the full support of $f_j$ is compact in
$\Omega_N$. Taking the limit proves membership in the slice's
$H_0^1$ space. This proof works for every choice of unfixed coordinates.
Dropping obstacles contributed by the fixed coordinates enlarges that
domain and preserves membership by zero extension.

The continuous map $x\mapsto(x,Y)$ sends a connected subset of $S_Y$
into one connected component of $\Omega_N$. The final assertion now
follows from Lemma~\ref{setup:ground}. In particular it holds whenever
two points can be joined by motion of the first coordinate entirely
inside $S_Y$.
\end{proof}

The slicing statement uses compactly supported approximation rather
than a boundary trace theorem. It therefore applies even when the
allowed slice has several components or an irregular boundary.
Moreover $m(Y)$ is finite for almost every bath, and baths outside the
allowed bath domain have $m(Y)=0$. Zero marginal slices contribute
nothing to the integral formulas below.

\subsection{The price of removing a part of each slice}

The triangle inequality in Equation~\eqref{setup:occupation} gives
$B(\Psi)\le B(\Phi)$. A lower bound for the modulus therefore needs a
separate phase estimate. We will retain points that communicate through
the open allowed slice and control the volume that was removed.

\begin{lemma}[Phase transfer]\label{setup:phase-transfer}
Let $\{(x,Y):x\in U_Y\}$ be measurable. Suppose that for almost every
bath with $m(Y)>0$, the set $U_Y$ is contained in one connected
component of $S_Y$, or is empty. Define
\[
 A_Y=\Lambda_K\setminus U_Y,\qquad
 \delta_{\rm geom}=\int m(Y)\frac{|A_Y|}{V}\,\dd Y.
\]
Then
\begin{equation}\label{setup:phase-loss}
 B(\Psi)\ge B(\Phi)-4\sqrt{\delta_{\rm geom}}.
\end{equation}
In particular, if $B(\Phi)\ge c_*>0$ and
$\delta_{\rm geom}\le c_*^2/64$, then $B(\Psi)\ge c_*/2$.
\end{lemma}

\begin{proof}
By Lemma~\ref{setup:slices}, the real part of
$\Psi(x,Y)\overline{\Psi(y,Y)}$ equals $\Phi(x,Y)\Phi(y,Y)$
whenever $x,y\in U_Y$. This includes the case in which the wave
function vanishes there. At all other points their difference is
between zero and $2\Phi(x,Y)\Phi(y,Y)$. Thus the double-integral
formula for the occupation gives
\[
 B(\Phi)-B(\Psi)
 \le\frac2V\int\!\!\int\!\!\int
   \bigl(\ind_{A_Y}(x)+\ind_{A_Y}(y)\bigr)
   \Phi(x,Y)\Phi(y,Y)\,\dd x\,\dd y\,\dd Y.
\]
For the term containing $\ind_{A_Y}(x)$, apply Cauchy--Schwarz to
\[
 V^{-1/2}\Phi(x,Y),\qquad
 V^{-1/2}\ind_{A_Y}(x)\Phi(y,Y).
\]
Their squared norms, with respect to $\dd x\,\dd y\,\dd Y$, are
$1$ and $\delta_{\rm geom}$, respectively. The corresponding term is
therefore at most $\sqrt{\delta_{\rm geom}}$. Interchanging $x$ and
$y$ proves the same bound for the other term, giving
Equation~\eqref{setup:phase-loss}. The final assertion follows by
substitution. This argument chooses no phase as a function of the bath;
all quantities integrated are already measurable.
\end{proof}

The rest of the proof must now supply two estimates with compatible
constants: a lower bound $B(\Phi)\ge c_*>0$ that remains valid as the
lower threshold on $D$ is increased, and a retained region for which
$\delta_{\rm geom}$ becomes sufficiently small. Lemma~\ref{setup:phase-transfer}
will then recover the occupation of every original ground vector.
